\section{Scripture, Rule of Faith, and Rival Christianities}

\subsection{Before the disputed text comes the apostolic Church}

\emph{On the Prescription of Heretics} is less a commentary than a challenge
to jurisdiction. A Roman legal prescription could stop a case before its
merits were tried. Tertullian likewise asks why teachers outside apostolic
communion may claim the Church's Scriptures and rearrange them against the rule
by which the Church received them. Chapters 19--21 move from Christ to apostles,
apostolic preaching, and churches founded by or joined to them. Chapter 32
challenges rival communities to display the succession and original of their
episcopal chairs. Chapter 36 points readers toward apostolic churches and
celebrates Rome's witness.

This is neither a modern archival proof nor the doctrine that no biblical
argument matters. Tertullian uses historical continuity, public teaching,
shared rule, and scriptural custody against schools he believes manufacture
private systems. When chapter 36 speaks of ``authentic writings'' read in
apostolic churches, the phrase concerns authoritative texts in their ecclesial
homes; it should not be inflated into a claim that first-century autograph
manuscripts still sat in each chair.

Chapter 13 states a rule of faith: one God, creator of the world through his
Word; the Word sent into the Virgin, born as Jesus Christ, crucified, raised,
taken into heaven, seated at the Father's right hand, sending the Holy Spirit,
and coming to receive the saints into eternal life and judge the wicked with
the restoration of the flesh. The rule is narrative and baptismal before it is
a catalog of technical terms. It gives Tertullian a boundary within which he
will argue fiercely and, in \emph{On the Veiling of Virgins} 1, distinguish
unchangeable faith from disciplines capable of development under the
Paraclete.

\subsection{Athens and Jerusalem in context}

Chapter 7 asks, ``What has Athens to do with Jerusalem?'' The maxim is often
made into Tertullian's rejection of all reason, learning, or philosophy. Its
immediate targets are philosophical genealogies assigned to heresies and a
curiosity that refuses the received Christ. The author making the attack uses
Stoic physics, dialectic, rhetorical classification, medicine, etymology, and
law across his corpus. The tension is real but more fruitful than the slogan:
he distrusts philosophy as an alternative authority while repeatedly making
its instruments serve ecclesial argument.

The same control applies to his confidence that searching has an end. The
prescription does not end all interpretation; Tertullian's books are evidence
of relentless interpretation. It ends a search imagined as permanent
instability after the apostolic rule has been found. His later appeal to new
prophetic discipline will test how fixed teaching and continuing divine speech
can coexist.

\subsection{Marcion and the archive of an opponent}

No rival occupied Tertullian longer than Marcion. Across five books
\emph{Against Marcion} defends the unity of the creator and the Father of
Jesus, the continuity of law, prophets, Gospel, and apostle, the goodness of
flesh, and the identity of Christ. Books 4 and 5 follow Marcion's gospel and
Pauline collection closely enough to make Tertullian an indispensable witness
to Marcionite Scripture.

Indispensable does not mean transparent. Tertullian quotes, paraphrases,
accuses, and arranges an opponent's text in order to defeat it. Modern
reconstruction must compare every hostile witness and textual form; his five
books are not a photographed Marcionite Bible. His repeated charge that Marcion
``mutilated'' Luke also assumes the priority and integrity of the ecclesial
form he knows. The polemic is evidence for competing textual histories, not the
final critical demonstration of their sequence.

The opening of book 1 makes the author's own revision history unusually visible.
Tertullian says an early work against Marcion was composed too hastily, that a
second form was lost to him when a later apostate fraudulently copied and
circulated it, and that the present work was expanded anew. The notice warns
against speaking of a Tertullianian ``book'' as if one edition left his desk and
never changed. Polemic, betrayal, authorial enlargement, copying, and later
manuscript loss all stand between argument and reader.

\subsection{A canon still visible in argument}

Tertullian's corpus supplies early evidence for Latin biblical reception, but
not one detached canon table. \emph{On the Apparel of Women} 1.3 defends Enoch
as useful to Christians. He acknowledges that it was not received in the
Jewish scriptural instrument known to him, then polemically explains that
rejection by the book's witness to Christ. His adversarial explanation is not a
neutral history of canon and cannot erase the diverse Jewish reception of
Enochic writings. \emph{On Modesty} 10 rejects the \emph{Shepherd} as an
adulterers' scripture, while chapter 20 attributes Hebrews to Barnabas and
reasons from it. \emph{On Baptism} 17 rejects the use made of the \emph{Acts of
Paul and Thecla}. These decisions arise within disputes over dress, absolution,
and ministry. They show boundaries being defended and contested, not a private
author issuing an ecumenical canon.

His citations are also major witnesses to the Old Latin Bible. Sometimes he
appears to quote a received Latin form; sometimes he renders Greek freely,
adapts wording to syntax or argument, or cites from memory. A modern reader
cannot identify ``Tertullian's Bible'' by back-translating every Latin sentence
as though it came from one codex. Critical editions and the Vetus Latina
apparatus are needed to distinguish textual witness from authorial handling.

\subsection{The adversarial manufacture of ``Jew''}

The work transmitted as \emph{Against the Jews} argues that the Gentiles have
received God's law and promises in Christ, reads Daniel's weeks, and treats
circumcision, Sabbath, sacrifice, and covenant typologically or
supersessionistically. Chapters 9--14 substantially overlap \emph{Against
Marcion} book 3, and the unity, completion, and authorship of that range remain
disputed. It should not be cited as an evenly secure single composition.

The rhetoric is historically consequential even where authorship is secure.
Tertullian makes ``the Jews'' a theological opponent within Christian claims to
Israel's Scripture. Polemical representation cannot substitute for description
of diverse Jewish people in Roman Africa or elsewhere. Recent scholarship,
including St\'ephanie Binder's study of Jews in Tertullian's literary
imagination, emphasizes that his portrayals shift with genre and target rather
than reproducing one social encounter. The Christian continuity he defends
with Israel's Scriptures must therefore be narrated alongside the harm made
possible when a rhetorical adversary becomes an undifferentiated people.

Tertullian's scriptural achievement and failure meet here. He refuses a Gospel
cut loose from creation, Israel, flesh, and public apostolic memory. Yet his
boundary-making can make living opponents exist only as exhibits inside the
Christian case. Source-first admiration cannot keep the first fact and hide the
second.
